% Generated draft: requires mathematical and source review.
\documentclass[11pt]{article}
\usepackage[margin=25mm]{geometry}
\usepackage{fontspec}
\setmainfont{Latin Modern Roman}
\usepackage{amsmath,amssymb,mathtools,unicode-math}
\setmathfont{Latin Modern Math}
\usepackage{enumitem,needspace,xurl,hyperref}
\setlength{\parindent}{0pt}
\setlength{\parskip}{5pt}
\setlength{\emergencystretch}{3em}
\newcommand{\N}{\mathbb{N}}
\newcommand{\Z}{\mathbb{Z}}
\newcommand{\Q}{\mathbb{Q}}
\newcommand{\R}{\mathbb{R}}
\newcommand{\C}{\mathbb{C}}
\newcommand{\F}{\mathbb{F}}
\newcommand{\A}{\mathbb{A}}
\newcommand{\PP}{\mathbb{P}}
\newcommand{\E}{\mathbb{E}}
\newcommand{\eps}{\varepsilon}
\newcommand{\dd}{\,\mathrm{d}}
\begin{document}
\section*{UnsolvedMath ID: 30003591}
\textbf{Source problem code:} catalogue entry
\subsection{Definitions and mathematical statement}
Write \(n=\dim X\), and let \(D\) denote the divisor intended in the source. Let \(k\) denote the ground field, which the supplied record does not specify. Whenever the sheaves \(\mathcal O_X(mD)\) are defined and their cohomology groups are finite-dimensional over \(k\), set
\[
\widehat h^i(D):=n!\limsup_{\substack{m\to\infty\\m\in\Z_{>0}}}\frac{\dim_k H^i(X,\mathcal O_X(mD))}{m^n}.
\]
Here \(i\) is a cohomological degree; the question concerns higher cohomology, hence positive degrees, although the source does not explicitly specify the range of \(i\).

The source asks two related questions:
\begin{enumerate}
\item Is there a natural convex-geometric construction, analogous to the Newton--Okounkov interpretation of divisor volume, which realizes \(\widehat h^i(D)\) as a Euclidean volume?
\item For every admissible \(X,D,i\), does the finite limit
\[
\lim_{\substack{m\to\infty\\m\in\Z_{>0}}}\frac{\dim_k H^i(X,\mathcal O_X(mD))}{m^n}
\]
exist, so that
\[
\widehat h^i(D)=n!\lim_{\substack{m\to\infty\\m\in\Z_{>0}}}\frac{\dim_k H^i(X,\mathcal O_X(mD))}{m^n}?
\]
\end{enumerate}
The quantifier ``every admissible \(X,D,i\)'' remains conditional: the supplied record does not state the admissible geometric class, ground-field hypotheses, or divisor hypotheses. In particular, projectivity, smoothness, and whether \(D\) is Cartier are not specified. If only certain positive multiples of \(D\) define the intended sheaves, the indexing sequence must also be specified; the displayed source formula uses all positive integers. The convex-geometric question likewise leaves the object, its dimension, its permitted auxiliary choices, and the normalization of Euclidean volume unspecified. These omissions prevent a fully determined universal formulation.
\subsection{Short English statement}
The problem asks for a natural convex-geometric volume interpretation of asymptotic higher cohomology and asks whether its defining \(\limsup\) is always an ordinary limit. The intended geometric hypotheses and the meaning of the proposed volume construction require clarification. The supplied registry labels the problem open; this draft does not independently verify that status.
\subsection{Sources}
\begin{itemize}
\item Supplied local UnsolvedMath registry record, permanent ID \(30003591\), source problem code \texttt{catalogue entry}, titled ``Convex Geometry and Limits of Asymptotic Higher Cohomology.'' The record gives proposed year 2017 and names no individual proposer.
\item Supplied report citation: \emph{Mini-Workshop: Positivity in Higher-dimensional Geometry: Higher-codimensional Cycles and Newton--Okounkov Bodies} (2018), \emph{Oberwolfach Reports}, pp.\ 2631--2657, DOI \texttt{10.4171/owr/2017/43}. Supplied source link: \url{https://doi.org/10.4171/owr/2017/43}. The background describes the report as documenting workshop discussions and open problems.
\item Supplied literature reference: ``Kuronya, asymptotic cohomological functions,'' \url{https://arxiv.org/abs/math/0501491}.
\item Supplied literature reference: ``Li, asymptotic cohomological functions and volumes,'' \url{https://arxiv.org/abs/2607.09021}.
\item Supplied literature reference: ``original OWR report,'' \url{https://doi.org/10.4171/OWR/2017/43}.
\item Literature assessment dated exactly \texttt{2026-08-21}: the supplied curation labels the problem open, attributes continuity and numerical invariance of asymptotic cohomological functions to Kuronya, states that the normalized higher-cohomology \(\limsup\) is not known to be a limit in general, and distinguishes the cited 2026 progress on volume limits from this higher-cohomology question. This is supplied dated curation derived from the OWR material, not an independently verified assessment; its time-sensitive status and cited claims require independent checking.
\end{itemize}
\end{document}
